% ============================================================================
% MUSTERLÖSUNG DS 4 - Gruppe A + B
% ============================================================================

\documentclass[11pt,a4paper]{article}

\usepackage[utf8]{inputenc}
\usepackage[T1]{fontenc}
\usepackage[ngerman]{babel}
\usepackage[left=2cm,right=2cm,top=1.8cm,bottom=1.5cm]{geometry}
\usepackage{helvet}
\usepackage{xcolor}
\usepackage{tabularx}
\usepackage{array}
\usepackage{enumitem}
\usepackage{tcolorbox}
\usepackage{fancyhdr}
\usepackage{lastpage}

\renewcommand{\familydefault}{\sfdefault}

\definecolor{spanisch}{HTML}{c41e3a}
\definecolor{grau}{HTML}{666666}
\definecolor{gruppeB}{HTML}{1565c0}
\definecolor{loesung}{HTML}{c62828}

\pagestyle{fancy}
\fancyhf{}
\fancyhead[L]{\small\textcolor{grau}{Spanisch Spätbeginner -- MUSTERLÖSUNG}}
\fancyhead[R]{\small\textcolor{grau}{ML-04}}
\fancyfoot[C]{\small\thepage/\pageref{LastPage}}
\renewcommand{\headrulewidth}{0.5pt}

\newcommand{\lsg}[1]{\textcolor{loesung}{\textbf{#1}}}

\newtcolorbox{gruppenbox}[1][]{
    colback=white,
    colframe=spanisch,
    fonttitle=\bfseries\color{white},
    colbacktitle=spanisch,
    title=#1,
    boxrule=1pt,
    arc=0pt,
    left=5pt, right=5pt, top=5pt, bottom=5pt
}

\newtcolorbox{gruppeBbox}[1][]{
    colback=white,
    colframe=gruppeB,
    fonttitle=\bfseries\color{white},
    colbacktitle=gruppeB,
    title=#1,
    boxrule=1pt,
    arc=0pt,
    left=5pt, right=5pt, top=5pt, bottom=5pt
}

\begin{document}

\begin{center}
{\Large\textbf{\textcolor{loesung}{MUSTERLÖSUNG}}}\\[0.3em]
{\large DS 4 -- Punto final / La vida laboral}
\end{center}

\vspace{0.5em}
\hrule
\vspace{1em}

% ============================================================================
% GRUPPE A
% ============================================================================
\begin{gruppenbox}[Gruppe A: Punto final + Evaluación (S. 33--34)]

\textbf{Aufgabe 1: E-Mail an einen Austauschschüler (10 P.) -- Beispiel}

\lsg{¡Hola Pablo!}

\lsg{Me llamo Laura y tengo 16 años. Vivo en Rostock, una ciudad en el norte de Alemania. Mi barrio se llama Südstadt y es muy tranquilo.}

\lsg{En mi barrio hay muchas tiendas, un parque grande y un polideportivo. La piscina está cerca de mi casa. Mi barrio es bonito porque hay muchos árboles.}

\lsg{Cuando vienes, podemos ir al cine o al centro comercial. También quiero ir contigo a la playa en Warnemünde.}

\lsg{¿Y tú? ¿Cómo es tu barrio? ¿Qué te gusta hacer?}

\lsg{¡Hasta pronto!}

\lsg{Laura}

\vspace{0.5em}

\textbf{Aufgabe 2: Evaluación ser/estar/hay (6 P.)}
\begin{enumerate}[label=\alph*), leftmargin=2em, nosep]
    \item Mi barrio \lsg{es} muy bonito.
    \item En el centro \lsg{hay} muchas tiendas.
    \item La piscina \lsg{está} cerca de mi casa.
    \item Nosotros \lsg{somos} de Madrid.
    \item ¿Dónde \lsg{está} el cine?
    \item En mi calle no \lsg{hay} parques.
\end{enumerate}

\vspace{0.5em}

\textbf{Aufgabe 3: Orte und Aktivitäten (5 P.)}

1 = \lsg{b}, 2 = \lsg{a}, 3 = \lsg{d}, 4 = \lsg{c}, 5 = \lsg{e}

\vspace{0.5em}

\textbf{Aufgabe 4: Konjugation von ir (6 P.)}
\begin{enumerate}[label=\alph*), leftmargin=2em, nosep]
    \item Yo \lsg{voy} al cine los sábados.
    \item ¿Adónde \lsg{vas} después de clase?
    \item Nosotros \lsg{vamos} a la discoteca.
    \item Mis amigos \lsg{van} al parque.
    \item María \lsg{va} a la biblioteca.
    \item ¿\lsg{Vais} al polideportivo?
\end{enumerate}

\vspace{0.5em}

\textbf{Aufgabe 5: Fragen für das Partnergespräch (8 P.) -- Beispiele}
\begin{enumerate}[leftmargin=2em, nosep]
    \item \lsg{¿Hay un parque en tu barrio?}
    \item \lsg{¿Cómo es tu barrio? ¿Es tranquilo o ruidoso?}
    \item \lsg{¿Adónde vas los fines de semana?}
    \item \lsg{¿Qué te gusta hacer en tu tiempo libre?}
\end{enumerate}

\end{gruppenbox}

\vspace{1em}

% ============================================================================
% GRUPPE B
% ============================================================================
\begin{gruppeBbox}[Gruppe B: La vida laboral (S. 106--107)]

\textbf{Kasten 1: Arbeitsleben DE/ES (4 P.)}

horario partido = \lsg{geteilte Arbeitszeit}

contratos temporales = \lsg{befristete Verträge}

desempleo juvenil / tasa de paro = \lsg{Jugendarbeitslosigkeit}

formación dual / aprendizaje = \lsg{duale Ausbildung}

\vspace{0.5em}

\textbf{Aufgabe 1: Zusammenfassung Arbeitsleben in Spanien (8 P.) -- Beispiel}

\lsg{En España, el horario laboral es diferente porque muchos españoles tienen una pausa larga al mediodía. Ya que hay mucho desempleo juvenil, muchos jóvenes tienen contratos temporales. Además, los españoles trabajan muchas horas, pero la productividad no es tan alta como en Alemania. Por eso, muchos jóvenes buscan trabajo en el extranjero.}

\vspace{0.5em}

\textbf{Aufgabe 2: Vergleich DE/ES (9 P.) -- Beispiele}
\begin{enumerate}[leftmargin=2em, nosep]
    \item \lsg{En Alemania hay horario continuo, mientras que en España muchos tienen horario partido con una pausa de 2-3 horas al mediodía.}
    \item \lsg{La tasa de desempleo juvenil en España es muy alta (más del 30\%), en cambio en Alemania es relativamente baja (menos del 10\%).}
    \item \lsg{En Alemania existe la formación dual, sin embargo en España la formación profesional es menos práctica y muchos jóvenes van a la universidad.}
\end{enumerate}

\vspace{0.5em}

\textbf{Aufgabe 3: Synonyme (5 P.)}

el empleo = \lsg{el trabajo, la ocupación}

la formación profesional = \lsg{el aprendizaje, la capacitación}

el desempleo = \lsg{el paro, estar sin trabajo}

las cualificaciones = \lsg{las competencias, las habilidades}

la jornada laboral = \lsg{el horario de trabajo, las horas de trabajo}

\vspace{0.5em}

\textbf{Aufgabe 4: Gliederung Kurzvortrag (10 P.) -- Beispiel}

\textbf{1. Einleitung:}

\lsg{Mi profesión ideal es ser médico/a. Me interesa este trabajo porque quiero ayudar a las personas.}

\textbf{2. Hauptteil:}

\lsg{Para ser médico necesito responsabilidad, empatía y buenas notas en biología y química. Tengo que estudiar medicina durante 6 años y después hacer prácticas en un hospital.}

\textbf{3. Schluss:}

\lsg{En conclusión, este trabajo es ideal para mí porque me gusta trabajar con personas y quiero tener un trabajo con sentido. Además, los médicos tienen buenas oportunidades de trabajo.}

\end{gruppeBbox}

\end{document}
